% commentaire.tex — f167r, vv. 380–422 (conversion initiale des travaux antérieurs)
\comm{381}{χρέμπτεται.}{Elle \italique{tousse}, et non elle crache, comme traduisent les dictionnaires. L'ancien français avait un terme excellent, qui s'est perdu avec tant d'autres : « Il se lève, compose sa robe, sa barbe et ses yeux à la modestie, et ayant craché et toussi clair, il commença ». D'Aubigné, \italique{Avent. du baron de Foeneste}, I. IV, ch. 10. (Note de Willems.)}
\comm{393}{τὰς οἰνοπίπας,}{Encore une fois, c'est le scholiaste, disposant de meilleurs manuscrits, qui nous permet de revenir à la leçon originale : après avoir évoqué Symmaque, il rapproche non sans raison \grec{οἰνοπίπης} de \grec{παιδοπίπης} (on aurait pu rajouter : \grec{γυναικοπίπης}, \grec{παρθενοπίπης} ou encore le \grec{πυροπίπης} / πυρροπίπης des \italique{Cavaliers}, v. 407) ; c'est-à-dire que le suffixe ne renverrait pas aux verbes \grec{πιπίζειν} ou \grec{πιπίσκειν} mais -\grec{οπίπης} serait de la famille de \grec{ὀπιπεύειν}, \italique{guigner}, \italique{guetter}, \italique{épier} (cf. aussi Hésychius (ο 1006) : \grec{ὀπῖπα· ἐξαπατᾷ, ἀπατεών, ἢ ἀπατῶν}). 
Aristophane ne va pas manquer de surenchérir : la difficile relation des femmes au vin sera allègrement mise en scène dans la suite de la pièce.}
\comm{394}{τὰς οὐδὲν ὑγιές.}{Érasme a compilé l'expression dans ses \italique{Adages} (n° 738, \latin{nihil sanum}).}
\comm{400}{τοὺς ἄνδρας ἡμῶν· ὥστ᾽ ἐάν τις πλέκῃ.}{Le vers 400 pèche contre le mètre. Willems y supplée en lisant <ξυμ>πλέκῃ ; il s'en justifie dans la \nieme{11} note, reproduite à la fin.}
\comm{400–1}{ὥστ᾽ ἐάν τις πλέκῃ / γυνὴ στέφανον, ἐρᾶν δοκεῖ·}{Le scholiaste s'est montré court sur ce travers qu'Euripide impute aux femmes tressandières : l'on ne sait de laquelle de ses pièces (perdues) le reproche est tiré.
Toutefois, tresser une couronne pour son amant(e) est chose bien attestée ; et le scholiaste de citer un vers de Sappho\nc{Citation corrompue ; le manuscrit de Ravenne donne : αυταόρα αἱστεφανηπλόκουν. La lecture proposée est de mon cru. Théodore Reinach, dubitatif, reprenait la conjecture de Hermann : αὐτὰρ ὀραῖαι στεφανηπλόκευν. Ce fragment de Sappho est le \nieme{70} de l'édition Reinach-Puech aux Belles Lettres (1937) et le \nieme{125} de l'édition d'Éva-Maria Voigt, \italique{Sappho et Alcaeus} (1971).} :
\citComm{νεωτέρων καὶ ἐρωτικῶν τὸ στεφανοπλοκεῖν. πρὸς τὸ ἔθος, ὅτι ἐστεφανηπλόκουν αἱ παλαιαί.\\
Σαπφώ· αὐτὰ κόραις στεφαναπλόκεο\\
Sappho : «~toi qui te faisais tresser des couronnes par des jeunes filles~».}
Athénée (XV, 670d) confirme la coutume d'offrir des couronnes à son être aimé :
\citComm{στεφανοῦσιν δὲ τὰς τῶν ἐρωμένων θύρας ἤτοι τιμῆς χάριν καθαπερεί τινος θεοῦ τὰ πρόθυρα [στεφανοῦσιν αὐτῶν], ἢ οὐ τοῖς ἐρωμένοις ἀλλὰ τῷ Ἔρωτι ποιούμενοι τὴν τῶν στεφάνων ἀνάθεσιν τοῦ μὲν Ἔρωτος τὸν ἐρώμενον ἄγαλμα, τούτου δὲ ναὸν ὄντα τὴν οἴκησιν στεφανοῦσι [τὰ τῶν ἐρωμένων πρόθυρα]: διὰ ταῦτα δὲ καὶ θύουσιν ἔνιοι ἐπὶ ταῖς τῶν ἐρωμένων θύραις.\\
Quant aux couronnes et aux guirlandes dont les amants ornent les portes des personnes qu'ils aiment, ils les y mettent ou pour leur rendre des honneurs, comme couronnant le porche d'une divinité ; ou, c'est à l'amour qu'ils les offrent en les posant, et non aux personnes qui sont aimées ; car ces personnes sont pour les amants les images mêmes du dieu d'amour; et la maison où elles demeurent en étant le temple, c'en est aussi le porche qu'ils ornent de couronnes, en les mettant aux portes de ces personnes. Voilà aussi pourquoi on immole quelquefois des victimes aux portes des personnes qu'on aime.}}
\comm{401–4}{κἂν ἐκβάλῃ ... τῷ Κορινθίῳ ξένῳ.}{\italique{L'étranger de Corinthe}, c'est Bellérophon, objet d'une coupable passion de la part de la princesse Sthénébée dans la pièce éponyme d'Euripide. Cette dernière, voyant ses avances repoussées par le chaste héros, l'accuse devant le roi son mari d'avoir essayé d'abuser d'elle. Le roi ordonne à Bellérophon de partir pour la Carie. Il y reçoit pour mission d'affronter la Chimère. C'est pendant le temps que Bellérophon est en Carie, et que Sthénébée le croit mort, qu'elle verse des sortes de libations à la mémoire du héros, comme l'on fait au jeu du cottabe avec les personnes aimées. Athénée (X, 427e) nous a conservé ce fragment\nc{fragment 3, \italique{Sthénébée}, de l'édition Jouan-van Looy, Belles Lettres, t. VIII, \nieme{3} partie $=$ fr. 664 de l'édition Kannicht $=$ fr. 883 de l'édition Mette. Traduction Jouan \& van Looy.} :
\citComm{πεσὸν δέ νιν λέληθεν οὐδὲν ἐκ χερός,\\
ἀλλ᾽ εὐθὺς αὐδᾷ· `τῷ Κορινθίῳ ξένῳ'\\
<La Nourrice>\\
\francais{Rien ne tombe, en lui échappant des mains,\\
qu'elle ne s'écrie aussitôt : « Pour l'hôte de Corinthe ! »}}Les maris athéniens prenaient leur femme pour autant de Sthénébées.\\
Ce passage est à rapprocher d'un fragment d'Aristophane, rapporté par Diogène Laërte\nc{\italique{Vies}, VIII, 1 (Pythagore), 34. Correspond au fragment 320(305) de l'édition Kassel-Austin (\italique{Poetae Comici Graeci}) des fragments d'Aristophane.} :
\citComm{Τὰ δὲ πεσόντα μὴ ἀναιρεῖσθαι, ὑπὲρ τοῦ ἐθίζεσθαι μὴ ἀκολάστως ἐσθίειν ἢ ὅτι ἐπὶ τελευτῇ τινος· καὶ Ἀριστοφάνης δὲ τῶν ἡρώων φησὶν εἶναι τὰ πίπτοντα, λέγων ἐν τοῖς Ἥρωσι,\\
Μηδὲ γεύεσθ' ἅττ' ἂν ἐντὸς τῆς τραπέζης καταπέσῃ.\\
$[$ Pythagore $]$ interdit de ramasser ce qui tombe de la table pour habituer à manger avec tempérance, ou bien encore parce que cela est destiné aux morts. Ce qui tombe de la table est pour les héros, suivant Aristophane ; car il dit dans les \italique{Héros} :
\francais{« Ne goûtez point à ce qui tombe de la table. »}}}
\comm{405–6}{εὐθὺς ἁδελφὸς λέγει· / `τὸ χρῶμα τοῦτό μ᾽ οὐκ ἀρέσκει τῆς κόρης.'}{Il la soupçonne d'être enceinte. La pièce d'Euripide à laquelle se réfère Aristophane n'est pas identifiée ; pourtant il semble clair qu'il s'agit d'une citation (ne serait-ce que par l'absence de résolution au vers 406).}
\comm{406–7}{γυνή τις ὑποβαλέσθαι βούλεται / ἀποροῦσα παίδων, οὐδὲ τοῦτ᾽ ἔστιν λαθεῖν.}{Le thème de la \italique{supposition d'enfants}\nc{Quand les femmes font passer pour leur des enfants – mâles le plus souvent ! – auprès de leur mari.} est récurrent dans les \italique{Thesmophories} (cf. vv. 340, 502-516, 564-565). À en croire Satyros dans sa \italique{Vita Euripidis}\nc{\italique{The Oxyrhynchus Papyri}, IX.1776, fr. 39, Col. vii, 1-24}, c'est à Euripide qu'on doit son omniprésence dans la Comédie Nouvelle... et le biographe d'Euripide de donner du grain à moudre à ses contempteurs :
\citComm{πρὸς γ[υ]ναῖκα καὶ πατρὶ πρὸς υ[ἱὸ]ν καὶ θερά[πον]τι πρὸς δ[εσ]πότην, ἢ τ[ὰ] κατὰ τὰς π[ερι]πετείας, β[ια]σμοὺς παρθένων, \textbf{ὑποβολὰς παιδίων}, ἀναγνωρισμοὺς διά τε δακτυλίων καὶ διὰ δεραίων· ταῦτα γάρ ἐστι δήπου τὰ συνέχοντα τα τὴν νεωτέραν κωμῳδίαν, ἃ πρὸς ἄκρον ἤγα[γ]εν Εὐρπίδης, Ὁμήρου [ὄ]ντος ἀρχῆς\\
\francais{$[$le mari$]$ contre sa femme, le père contre son fils, le serviteur contre son maître ; ou lors des péripéties : violences de jeunes filles, \textbf{substitutions d'enfants}, reconnaissances par des anneaux ou des colliers ; hé bien c'est tous ces éléments qui composent la trame de la Nouvelle Comédie, éléments qu'Euripide avait conduit au sommet, Homère en étant le point de départ}}La tradition ne nous a rapporté qu'une seule supposition d'enfants chez Euripide : dans les \italique{Phéniciennes} (v. 28-31) :
\citComm{Πολύβου δέ νιν λαβόντες ἱπποβουκόλοι\\
φέρουσ᾽ ἐς οἴκους ἔς τε δεσποίνης χέρας\\
ἔθηκαν. ἣ δὲ τὸν ἐμὸν ὠδίνων πόνον\\
μαστοῖς ὑφεῖτο καὶ πόσιν πείθει τεκεῖν,}ainsi qu'une allusion dans l'\italique{Alceste} (v. 639) : Admète, désapprouvant l'attitude de son père qui a refusé de donner sa vie pour lui, suppose qu'il a été supposé.}
\comm{413}{δέσποινα γὰρ γέροντι νυμφίῳ γυνή.}{Stobée (IV, 22a109) nous a conservé la fameuse citation d'Euripide ; elle est tirée du \italique{Phoenix}\nc{fr. 3, \italique{Phoenix} de l'édition Jouan-van Looy aux Belles-lettres, tome VIII, \nieme{3} partie. Je reprends leur traduction. Ce fragment correspond au n°804 de l'édition Kannicht et 1139 de celle de Mette.} :
\citComm{μοχθηρόν ἐστιν ἀνδρὶ πρεσβύτῃ τέκνα\\
δίδωσιν†, ὅστις οὐκέθ᾽ ὡραῖος γαμεῖ·\\
δέςποινα γὰρ γέροντι νυμφίῳ γυνή\\
\francais{C'est un tracas pour un vieillard <qui est doté d'enfants> de se marier alors qu'il n'en a plus l'âge ; avec un vieil époux, une femme est la maîtresse.}}\italique{Quand on est jeune, on aime en fou ; Quand on est vieux, qui aime est fou.}}
\comm{416–7}{Μολοττικοὺς / τρέφουσι μορμολυκεῖα τοῖς μοιχοῖς κύνας.}{C'est le conseil et le vœu d'Hippolyte dans la pièce d'Euripide (vv. 645 et suiv. ; note de Willems) : 
\citComm{χρῆν δ᾽ ἐς γυναῖκα πρόσπολον μὲν οὐ περᾶν,\\
ἄφθογγα δ᾽ αὐταῖς συγκατοικίζειν δάκη\\
θηρῶν, ἵν᾽ εἶχον μήτε προσφωνεῖν τινα\\
μήτ᾽ ἐξ ἐκείνων φθέγμα δέξασθαι πάλιν.\\
\francais{Il faudrait que les femmes n'eussent point auprès d'elles de servantes,\\
mais qu'elles fussent servies par de muets animaux,\\
pour qu'elles n'eussent personne à qui parler,\\
ni qui pût à son tour leur adresser la parole.}}}
\comm{418–20}{ἃ δ᾽ ἦν ἡμῖν πρὸ τοῦ / αὐταῖς ταμιεύεσθαι προαιρούσαις λαβεῖν / ἄλφιτον ἔλαιον οἶνον.}{Je renvoie à la courte mais incisive note (la \nieme{12}) de Willems reproduite à la fin : il y fait un sort à la « déplorable conjecture de Scaliger » (qui avait substitué λαθεῖν à λαβεῖν au vers 419, donnant au passage un sens qui fait contre-sens).}
